\documentclass[12pt,a4paper]{article}

% ── Packages ──────────────────────────────────────────────────────────────────
\usepackage[margin=2.5cm]{geometry}
\usepackage{graphicx}
\usepackage{amsmath}
\usepackage{booktabs}
\usepackage{hyperref}
\usepackage{xcolor}
\usepackage{longtable}
\usepackage{array}
\usepackage{enumitem}
\usepackage{caption}
\usepackage[numbers,sort&compress]{natbib}
\usepackage{microtype}

\definecolor{argblue}{RGB}{26,79,138}
\definecolor{apsorange}{RGB}{232,119,34}

\hypersetup{
  colorlinks=true,
  linkcolor=argblue,
  citecolor=argblue,
  urlcolor=argblue,
}

% Renumber tables and figures with S prefix
\renewcommand{\thetable}{S\arabic{table}}
\renewcommand{\thefigure}{S\arabic{figure}}

% ── Title block ────────────────────────────────────────────────────────────────
\title{\textbf{Supporting Information}\\[6pt]
\large Query-GSAS: A Local Retrieval-Augmented Generation System
for Intelligent Access to GSAS-II Documentation}

\author{
  Pawan K. Tripathi \and
  Mathew J.\ Cherukara \and
  Brian H.\ Toby
}

\date{
  X-ray Science Division, Advanced Photon Source,
  Argonne National Laboratory, Lemont, IL 60439, USA
}

% ── Document ───────────────────────────────────────────────────────────────────
\begin{document}

\maketitle

\noindent This Supporting Information accompanies the manuscript
\textit{Query-GSAS: A Local Retrieval-Augmented Generation System for
Intelligent Access to GSAS-II Documentation} submitted to the
\textit{Journal of Applied Crystallography}.
It provides (i) detailed corpus and index statistics, (ii) a description
of the embedding models compared, (iii) the complete 40-question evaluation
set with per-question results, and (iv) an extended failure analysis.

\tableofcontents
\clearpage

% ─────────────────────────────────────────────────────────────────────────────
\section{Corpus and index statistics}
\label{sec:corpus}
% ─────────────────────────────────────────────────────────────────────────────

The \textit{Query-GSAS} documentation index is built from five source
categories.
Table~\ref{tab:corpus} reports verified statistics from a live ingest
run performed on 2026-07-31, using the
\texttt{BAAI/bge-base-en-v1.5} embedding model and the full index
(including the optional \textit{Powder Diffraction Crystallography}
textbook).

\begin{table}[htbp]
  \centering
  \caption{Verified source statistics for the \textit{Query-GSAS} documentation
           index.
           Word counts are approximate (whitespace-delimited tokens after HTML
           stripping).
           Chunk counts use a maximum chunk size of 1{,}200 characters with
           sentence-boundary splitting.
           Tutorial URLs are resolved dynamically from the GSAS-II
           \texttt{tutorialIndex.py} at ingest time; book URLs are resolved by
           probing the GitHub Pages site for chapter and section HTML files.}
  \label{tab:corpus}
  \begin{tabular}{lrrr}
    \toprule
    Source category & HTML pages & $\approx$Words & Chunks \\
    \midrule
    Home pages (installation, install guides) & 22  & 22{,}362  & 329   \\
    Help manual sections                       & 42  & 54{,}089  & 554   \\
    Tutorials (dynamically fetched)            & 65  & 166{,}694 & 1{,}437 \\
    Programmer's Guide (ReadTheDocs)           & 23  & 143{,}486 & 2{,}740 \\
    \textit{Powder Cryst.\ textbook (opt.)}    & 179 & 115{,}079 & 778   \\
    \midrule
    \textbf{Default index (without book)}      & \textbf{152} & \textbf{386{,}631} & \textbf{5{,}060} \\
    \textbf{Full index (with book)}            & \textbf{331} & \textbf{501{,}710} & \textbf{5{,}838} \\
    \bottomrule
  \end{tabular}
\end{table}

% ─────────────────────────────────────────────────────────────────────────────
\section{Embedding models}
\label{sec:embed}
% ─────────────────────────────────────────────────────────────────────────────

Two sentence embedding models were evaluated.
Table~\ref{tab:models} summarises their key properties.
All inference is performed locally via the ONNX runtime;
no data are transmitted to external servers.

\begin{table}[htbp]
  \centering
  \caption{Embedding models compared in the retrieval evaluation.
           Both run locally via the ONNX runtime without requiring PyTorch.
           \texttt{all-MiniLM-L6-v2} is accessed through ChromaDB's built-in
           \texttt{DefaultEmbeddingFunction};
           \texttt{BAAI/bge-base-en-v1.5} is accessed via the
           \texttt{fastembed} library.}
  \label{tab:models}
  \begin{tabular}{llrrll}
    \toprule
    Model & Shorthand & Dim & Size & Runtime & Library \\
    \midrule
    \texttt{all-MiniLM-L6-v2}     & MiniLM   & 384 & $\approx$90\,MB  & ONNX & chromadb \\
    \texttt{BAAI/bge-base-en-v1.5} & bge-base & 768 & $\approx$210\,MB & ONNX & fastembed \\
    \bottomrule
  \end{tabular}
\end{table}

The \texttt{bge-base-en-v1.5} model was trained on a large-scale English
corpus with contrastive learning objectives specifically targeting
retrieval tasks~\cite{bge2023embedding}.
Its 768-dimensional embeddings provide finer-grained semantic
discrimination compared with the 384-dimensional \texttt{all-MiniLM-L6-v2}
baseline.

\subsection{Retrieval configuration}
\label{sec:retrieval}

The production retrieval pipeline uses the following steps:

\begin{enumerate}[leftmargin=*]
  \item \textbf{Over-fetch}: retrieve the top-30 chunks by cosine similarity to
        the embedded query.
  \item \textbf{URL-diversity deduplication}: for each unique source URL, retain
        only the highest-scoring chunk.
        This prevents a single high-chunk-count page from occupying multiple
        result slots.
        For example, the tutorials index page (\texttt{tutorials.html}) contains
        59 chunks; without deduplication it would frequently dominate the result
        set for broad queries.
  \item \textbf{Selection}: return the top-6 chunks after deduplication as the
        context for LLM synthesis.
\end{enumerate}

% ─────────────────────────────────────────────────────────────────────────────
\section{Evaluation set}
\label{sec:eval}
% ─────────────────────────────────────────────────────────────────────────────

The evaluation set comprises 40 questions spanning six GSAS-II workflow
categories (Table~\ref{tab:categories}).
Questions were authored by domain experts and cross-checked against the
GSAS-II tutorials, help manual, and Programmer's Guide.
Ground-truth source URLs were identified by manually locating the
documentation page that most directly answers each question.
For questions with conceptual content also covered by the
\textit{Powder Diffraction Crystallography} textbook, the relevant
textbook section URL was added as an additional valid answer.

\begin{table}[htbp]
  \centering
  \caption{Evaluation question categories.}
  \label{tab:categories}
  \begin{tabular}{lcp{9cm}}
    \toprule
    Category & $n$ & Representative question \\
    \midrule
    Rietveld    & 10 & How do I start a Rietveld refinement for laboratory X-ray powder data? \\
    Sequential  & 8  & How do I set up a sequential Rietveld refinement for a temperature series? \\
    Structure   & 6  & How do I use charge flipping for structure solution in GSAS-II? \\
    Calibration & 6  & How do I calibrate an area detector geometry in GSAS-II? \\
    Scripting   & 5  & How do I write a Python script to run a GSAS-II Rietveld refinement? \\
    Export      & 5  & How do I create a CIF file for publication from a GSAS-II refinement? \\
    \midrule
    Total       & 40 & \\
    \bottomrule
  \end{tabular}
\end{table}

The full question set with ground-truth URLs is provided in
\texttt{evaluation/questions.json} in the repository.

% ─────────────────────────────────────────────────────────────────────────────
\section{Retrieval results}
\label{sec:results}
% ─────────────────────────────────────────────────────────────────────────────

\subsection{Overall metrics}
\label{sec:overall}

Table~\ref{tab:configs} compares retrieval performance across four
configurations.
Results were obtained by running the evaluation set against each
configuration using the same scorer (exact URL match or prefix match
against the ground-truth URL set).

\begin{table}[htbp]
  \centering
  \caption{Retrieval metrics for four pipeline configurations on the
           40-question evaluation set.
           R@$k$ = Recall at $k$ (proportion of questions where a correct
           source appears in the top $k$ results).
           MRR = Mean Reciprocal Rank.
           ``Book GT'' = book section URLs accepted as valid answers for
           questions whose content is covered by the textbook.}
  \label{tab:configs}
  \begin{tabular}{lcrrrr}
    \toprule
    Configuration & Chunks & R@1 & R@3 & R@6 & MRR \\
    \midrule
    MiniLM, no dedup, no book         & 5{,}055 & 0.325 & 0.600 & 0.725 & 0.464 \\
    MiniLM + URL dedup, no book       & 5{,}055 & 0.325 & 0.600 & 0.750 & 0.464 \\
    bge-base + URL dedup + book       & 5{,}838 & 0.325 & 0.675 & 0.825 & 0.493 \\
    \textbf{bge-base + URL dedup + book (book GT)} & \textbf{5{,}838} & \textbf{0.500} & \textbf{0.825} & \textbf{0.975} & \textbf{0.661} \\
    \bottomrule
  \end{tabular}
\end{table}

\subsection{Per-category breakdown}
\label{sec:percategory}

Table~\ref{tab:percategory} reports all retrieval metrics broken down by
question category for the final (bge-base, with book, book GT accepted)
configuration.

\begin{table}[htbp]
  \centering
  \caption{Per-category retrieval metrics for the \texttt{bge-base-en-v1.5}
           model with URL-diversity deduplication and the full index
           (5{,}838 chunks).
           Book section URLs are accepted as valid answers where applicable.}
  \label{tab:percategory}
  \begin{tabular}{lcrrrr}
    \toprule
    Category & $n$ & R@1 & R@3 & R@6 & MRR \\
    \midrule
    All         & 40 & 0.500 & 0.825 & 0.975 & 0.661 \\
    \midrule
    Rietveld    & 10 & 0.400 & 0.700 & 1.000 & 0.553 \\
    Sequential  & 8  & 0.375 & 0.750 & 1.000 & 0.608 \\
    Structure   & 6  & 0.333 & 1.000 & 1.000 & 0.639 \\
    Calibration & 6  & 1.000 & 1.000 & 1.000 & 1.000 \\
    Scripting   & 5  & 0.400 & 0.800 & 0.800 & 0.533 \\
    Export      & 5  & 0.600 & 0.800 & 1.000 & 0.707 \\
    \bottomrule
  \end{tabular}
\end{table}

\subsection{Per-question results}
\label{sec:perquestion}

Table~\ref{tab:perquestion} gives the complete per-question results for the
bge-base-en-v1.5 configuration with the full index and book ground-truth
expansion.
Columns show the question ID, category, Recall@1/3/6 (Y = hit, N = miss),
reciprocal rank (RR), and the highest-ranked retrieved URL
(abbreviated; full URLs in \texttt{evaluation/eval\_results\_bge.json}).

\begin{longtable}{clcccc>{\raggedright\footnotesize\arraybackslash}p{5.8cm}}
  \caption{Per-question retrieval results (\texttt{bge-base-en-v1.5}, full
           index, book GT accepted).
           URL prefixes: \texttt{tut/} =
           \texttt{.../GSAS-II-tutorials/};
           \texttt{book/} = \texttt{.../PowderCrystallography/};
           \texttt{rtd/} = \texttt{gsas-ii.readthedocs.io/en/latest/}.
           Column shows the rank-1 retrieved URL (may differ from ground
           truth for hits at $k>1$).}
  \label{tab:perquestion} \\
  \toprule
  Q & Category & @1 & @3 & @6 & RR & Top-1 retrieved URL \\
  \midrule
  \endfirsthead
  \multicolumn{7}{l}{\small\textit{(Table~\ref{tab:perquestion} continued)}} \\
  \toprule
  Q & Category & @1 & @3 & @6 & RR & Top-1 retrieved URL \\
  \midrule
  \endhead
  \midrule
  \multicolumn{7}{r}{\small\textit{(continued on next page)}} \\
  \endfoot
  \bottomrule
  \endlastfoot
  % ── Rietveld ──
   1 & Rietveld    & N & N & Y & 0.17 & tut/tutorials.html \\
   2 & Rietveld    & Y & Y & Y & 1.00 & book/HTML-templatese54.html \\
   3 & Rietveld    & N & N & Y & 0.17 & rtd/GSASIIscriptable.html \\
   4 & Rietveld    & Y & Y & Y & 1.00 & book/HTML-templatese76.html \\
   5 & Rietveld    & N & Y & Y & 0.33 & tut/tutorials.html \\
   6 & Rietveld    & N & Y & Y & 0.33 & rtd/graphics.html \\
   7 & Rietveld    & Y & Y & Y & 1.00 & book/HTML-templatese64.html \\
   8 & Rietveld    & Y & Y & Y & 1.00 & tut/LeBail/LeBailSucrose.htm \\
   9 & Rietveld    & N & Y & Y & 0.33 & tut/TOF Charge Flipping/ \\
  10 & Rietveld    & N & N & Y & 0.20 & tut/tutorials.html \\
  % ── Sequential ──
  11 & Sequential  & N & N & Y & 0.20 & tut/RMCProfile-I/ \\
  12 & Sequential  & N & N & Y & 0.17 & tut/tutorials.html \\
  13 & Sequential  & Y & Y & Y & 1.00 & tut/SeqRefine/SequentialTutorial.htm \\
  14 & Sequential  & Y & Y & Y & 1.00 & tut/SeqParametric/ParametricFitting.htm \\
  15 & Sequential  & Y & Y & Y & 1.00 & tut/SeqRefine/SequentialTutorial.htm \\
  16 & Sequential  & N & Y & Y & 0.50 & tut/tutorials.html \\
  17 & Sequential  & N & Y & Y & 0.50 & tut/tutorials.html \\
  18 & Sequential  & N & Y & Y & 0.50 & tut/TOF-CW Joint Refinement/ \\
  % ── Structure ──
  19 & Structure   & N & Y & Y & 0.33 & tut/CFXraySingleCrystal/ \\
  20 & Structure   & N & Y & Y & 0.50 & tut/k\_vec\_isodistort/ \\
  21 & Structure   & N & Y & Y & 0.50 & tut/k\_vec\_tutorial\_non\_zero/ \\
  22 & Structure   & Y & Y & Y & 1.00 & tut/CFXraySingleCrystal/ \\
  23 & Structure   & N & Y & Y & 0.50 & tut/StackingFaults-III/ \\
  24 & Structure   & Y & Y & Y & 1.00 & tut/MerohedralTwins/ \\
  % ── Calibration ──
  25 & Calibration & Y & Y & Y & 1.00 & tut/CalibrationTutorial/ \\
  26 & Calibration & Y & Y & Y & 1.00 & tut/DeterminingWavelength/ \\
  27 & Calibration & Y & Y & Y & 1.00 & book/HTML-templatese85.html \\
  28 & Calibration & Y & Y & Y & 1.00 & tut/FPAfit/FPAfit.htm \\
  29 & Calibration & Y & Y & Y & 1.00 & book/HTML-templatech11.html \\
  30 & Calibration & Y & Y & Y & 1.00 & tut/CalibrationTutorial/ \\
  % ── Scripting ──
  31 & Scripting   & N & Y & Y & 0.33 & tut/RMCProfile-III/ \\
  32 & Scripting   & Y & Y & Y & 1.00 & rtd/GSASIIscriptable.html \\
  33 & Scripting   & N & N & N & 0.00 & tut/CFsucrose/ \\
  34 & Scripting   & N & Y & Y & 0.33 & rtd/GSASIIscriptable.html \\
  35 & Scripting   & Y & Y & Y & 1.00 & tut/ParameterLimits/ \\
  % ── Export ──
  36 & Export      & Y & Y & Y & 1.00 & tut/CIFtutorial/ \\
  37 & Export      & Y & Y & Y & 1.00 & book/HTML-templatech15.html \\
  38 & Export      & Y & Y & Y & 1.00 & tut/RietPlot/PublicationPlot.htm \\
  39 & Export      & N & Y & Y & 0.33 & rtd/GSASIIscriptable.html \\
  40 & Export      & N & N & Y & 0.20 & tut/StackingFaults-I/ \\
\end{longtable}

% ─────────────────────────────────────────────────────────────────────────────
\section{Failure analysis}
\label{sec:failure}
% ─────────────────────────────────────────────────────────────────────────────

Of the 40 questions, 39 have at least one correct source in the top-6
results (Recall@6 = 0.975).
The single complete miss (Q33) and the patterns underlying
sub-optimal ranking at lower $k$ are discussed below.

\subsection{Complete miss: Q33}

\noindent\textbf{Question:} ``How do I access and read back refinement results
(lattice parameters, Rwp) from a GSAS-II script?''

\noindent\textbf{Ground truth:} \texttt{PythonScript/Scripting.htm},
\texttt{rtd/GSASIIscripts.html}

\noindent\textbf{Retrieved at rank 1--6:}
charge-flipping sucrose, Laboratory~X, Charge Flipping in GSAS.htm,
RMCProfile-I, MC simulated annealing, RMCProfile-IV.

\noindent\textbf{Root cause:}
The terms ``lattice parameters'' and ``Rwp'' appear as refinement outputs
throughout the GSAS-II tutorial corpus, including in unrelated tutorials
(charge flipping, RMC).
The scripting-specific vocabulary (\texttt{G2Project}, \texttt{get\_LastFitResults},
\texttt{histogram dict}) is absent from the question text.
The bge-base-en-v1.5 model cannot recover the scripting context
when the question uses only quantity names rather than API vocabulary.

\noindent\textbf{Potential remedies:}
\begin{itemize}[leftmargin=*]
  \item Add dedicated scripting-API documentation that associates
        \texttt{get\_LastFitResults()} with ``Rwp'' and ``lattice parameters''
        so that the relevant chunk is uniquely associated with the scripting
        source.
  \item Reformulate the question with scripting-specific terms (\emph{e.g.}\
        ``G2Project object'', ``get\_LastFitResults'').
  \item Implement hybrid retrieval (dense embedding + BM25 keyword scoring)
        to penalise non-scripting pages that contain these terms incidentally.
\end{itemize}

\subsection{Systematic low-rank patterns}

Several questions are correctly answered at Recall@6 but rank the
ground-truth page at position 4--6, reducing Recall@1 and MRR.
Two structural causes account for most cases.

\paragraph{Tutorials index page (\texttt{tutorials.html}) at rank 1.}
This single page has 59 indexed chunks and describes every tutorial in
one place.
For broad category-level questions (\emph{e.g.}\ ``How do I start a
Rietveld refinement?'' or ``How do I set up a sequential refinement?''),
\texttt{tutorials.html} ranks first because its summary vocabulary matches
many queries simultaneously.
URL-diversity deduplication already prevents this page from occupying
more than one of the six result slots, but it cannot prevent it from
taking rank~1.
Affected questions: Q1, Q5, Q10, Q12, Q16, Q17.

\paragraph{Closely related tutorial returned instead of target.}
For questions with specific tutorial targets, semantically similar
tutorials occasionally rank above the intended page.
Example: Q9 (atom positions and Uiso) retrieves the TOF charge-flipping
tutorial at rank~1 because both involve crystal structure refinement
with similar vocabulary.
Example: Q11 (sequential refinement setup) retrieves RMCProfile-I at
rank~1 because both describe importing multiple datasets.
In all such cases the correct page appears by rank~6, and the correct
structural tutorial type (sequential, Rietveld, \emph{etc.}) is
always represented among the six returned sources.

\subsection{Categories with perfect Recall@6}

\paragraph{Calibration (R@6 = 1.00, MRR = 1.00).}
Calibration tutorial names (\texttt{CalibrationTutorial},
\texttt{DeterminingWavelength}, \texttt{FPAfit}) are semantically
specific and not confused with other content categories.
All six calibration questions have the correct page at rank~1.

\paragraph{Structure (R@6 = 1.00).}
Structure-solution tutorials (charge flipping variants, merohedral twins,
stacking faults, k-vector search) have highly distinctive vocabulary.
The model reliably distinguishes them despite multiple charge-flipping
tutorials sharing ``charge flipping'' as a keyword.

\paragraph{Sequential and Rietveld (R@6 = 1.00).}
All questions are answered by rank~6 despite some requiring the correct
page to be recovered from behind \texttt{tutorials.html} or a related
tutorial.

\subsection{Model impact}

Switching from \texttt{all-MiniLM-L6-v2} (384-dim) to
\texttt{BAAI/bge-base-en-v1.5} (768-dim) with URL-deduplication, on the
default index (no book), raised R@6 from 0.725 to 0.825 (+10 percentage
points).
Table~\ref{tab:modelimpact} shows the per-category breakdown.

\begin{table}[htbp]
  \centering
  \caption{Per-category Recall@6 comparison between the MiniLM-L6-v2
           baseline (no URL dedup, no book) and the upgraded
           bge-base-en-v1.5 pipeline (URL dedup, with book and book GT).}
  \label{tab:modelimpact}
  \begin{tabular}{lrrr}
    \toprule
    Category & R@6 MiniLM & R@6 bge-base & $\Delta$ (pp) \\
    \midrule
    All         & 0.725 & 0.975 & $+$25.0 \\
    Rietveld    & 0.800 & 1.000 & $+$20.0 \\
    Sequential  & 0.500 & 1.000 & $+$50.0 \\
    Structure   & 0.833 & 1.000 & $+$16.7 \\
    Calibration & 0.667 & 1.000 & $+$33.3 \\
    Scripting   & 1.000 & 0.800 & $-$20.0 \\
    Export      & 0.600 & 1.000 & $+$40.0 \\
    \bottomrule
  \end{tabular}
\end{table}

The Scripting regression ($-$20 pp) is attributable to Q33 alone, which
is a structural vocabulary failure that persists regardless of embedding
model.
All other categories show improvements or no change with bge-base.
Note that the bge-base evaluation includes the textbook in both the index
and the ground truth; part of the improvement over MiniLM therefore
reflects the additional coverage from 778 textbook chunks, not solely the
embedding model upgrade.

\bibliographystyle{unsrtnat}
\bibliography{references}

\end{document}
